\chapterimage{tubes.jpg}

\chapter{Capa 4: Transporte}\label{sec:layer4}

\begin{minipage}{0.4\linewidth}
\begin{center}
\begin{bytefield}{16}
\bitbox{16}{Layer 7: Application} \\
\bitbox{16}{\color{color1} Layer 4: Transport} \\
\bitbox{16}{Layer 3: Internet} \\
\bitbox{16}{Layer 2: Network (LAN)} \\
\bitbox{16}{Layer 1: Physical} \\
\end{bytefield}
\end{center}
\end{minipage}
\begin{minipage}{0.6\linewidth}
\begin{center}
\includegraphics[width=\linewidth]{network_layer4.pdf}
\end{center}
\end{minipage}

\vspace{-0.75cm}

\subsection*{Capacidades}

La capa de \conceptRef{transport}{transporte} comunica procesos (programas)
que se ejecutan en \conceptRef{host}{hosts} (equipos conectados a la red) potencialmente distintos.
Estos procesos se identifican mediante direcciones de capa~4 llamadas \conceptRef{port}{puertos}.

La capa de transporte ofrece dos formas distintas de comunicación: sin conexiones (UDP)
y con conexiones (TCP).

UDP no establece una conexión con el host de destino porque la prioridad es
enviar unos datos de payload (carga útil) lo más rápido posible.
En su lugar, UDP envía paquetes llamados \conceptRef{user datagram}{datagramas de usuario} directamente al destino,
sin esperar ningún tipo de confirmación de recepción por parte de la capa UDP del destino.
Las cabeceras UDP son simples y pequeñas, de modo que el overhead (sobrecarga) de este protocolo también se mantiene pequeño.

TCP establece una conexión antes de enviar cualquier dato de payload
para asegurarse de que ambos extremos están listos. Además, TCP espera un \conceptRef{ACK}{acuse de recibo} (confirmación)
de la recepción de cada byte de datos de payload. En caso contrario, los datos se \conceptRef{retransmission}{retransmiten}
tras un periodo de espera, lo que le da a TCP su \conceptRef{reliability}{fiabilidad}. TCP también usa una forma muy ingeniosa
de limitar la velocidad de intercambio de datos (\conceptRef{flow control}{control de flujo}) para que ordenadores rápidos y lentos puedan coexistir
y comunicarse mediante TCP/IP. Para ofrecer todas estas prestaciones, TCP es un protocolo más complejo que UDP,
usa cabeceras más grandes e introduce un mayor overhead.

\begin{exercise}
¿Cuáles de las siguientes aplicaciones crees que usan UDP y TCP?
\begin{itemize}
\item Juegos multijugador
\item La WWW
\item BitTorrent
\item Streaming de vídeo
\end{itemize}
\end{exercise}

\subsection*{Protocolos}

\begin{center}
\includegraphics[width=0.95\linewidth]{protocols_layer4.pdf}
\end{center}

La capa de transporte está en lo más alto de la pila que proporciona el sistema operativo.
Prácticamente todas las aplicaciones y, por tanto, sus protocolos de comunicación,
emplean TCP, UDP (o incluso ambos) para intercambiar datos por la red.
En la figura anterior se muestran algunos protocolos importantes que usan TCP o UDP.

Tanto TCP como UDP encapsulan sus PDU en datagramas IP.
Generalmente se usa IPv4, pero IPv6 también es una posibilidad cuando está disponible.

\begin{remark}
Ni TCP ni UDP cambian en absoluto cuando se usa IPv6 en lugar de IPv4, \ie,
no existe un protocolo TCPv6. En su lugar, \inlineCode{tcp6} y otras expresiones que incluyen
TCP y 6 se refieren a TCP (el de siempre) sobre IPv6.
\end{remark}


\section{Direccionamiento de capa 4}\label{sec:layer4:addressing}

Las ``direcciones'' de TCP y UDP se llaman puertos y son enteros sin signo de $16$ bits.
Para las personas, los puertos se expresan principalmente de dos formas:
\begin{itemize}
\item Como un número, \eg\ 22, 80 o 443.
\item Como el nombre del servicio que suele usar ese puerto, \eg\ \concept{sshd}, \conceptRef{HTTP}{http}, \conceptRef{HTTPs}{https}.
\end{itemize}

Cuando se recibe una PDU de capa~4, el sistema operativo usa el puerto de su cabecera para saber cuál de los procesos
en ejecución debe recibir los datos.
A su vez, los procesos tienen que hacer \concept{bind} (asociarse) a un puerto TCP o a un puerto TCP (o a ambos)
para que el sistema operativo sepa de qué puerto esperan datos.

Solo un proceso puede hacer bind a un puerto a la vez para enviar datos, pero los puertos TCP y UDP
se tratan de forma independiente:
\begin{itemize}
 \item Si un proceso está asociado al puerto \otherBase{TCP/100}, no recibirá tráfico UDP dirigido al puerto 100.
 \item Si un proceso está asociado al puerto \otherBase{TCP/100}, otro proceso (o incluso el mismo) puede
   hacer bind a \otherBase{UDP/100}.  
\end{itemize}

El puerto número \otherBase{0} está reservado, así que el rango de puertos utilizable es \otherBase{1}--\otherBase{65535}.
Los puertos \otherBase{1}--\otherBase{1023} solo se asocian a
procesos con privilegios (\eg\ un servidor http), mientras que las aplicaciones cliente (\eg\ un cliente http)
solo pueden hacer bind a los puertos \otherBase{1024}--\otherBase{65535} en función de su disponibilidad.

\begin{exercise}
¿Cuántos servidores distintos se pueden ejecutar en una máquina para ofrecer alguna funcionalidad a otros hosts de la red?
\end{exercise}


\section{UDP -- User Datagram Protocol}\label{sec:layer4:udp}
\subsection{Formato del paquete}\label{sec:layer4:udp:packet_format}

Todos los \conceptRef{user datagram}{datagramas de usuario} que usa UDP tienen una cabecera de $8$ bytes y
la siguiente estructura:

\begin{center}
\begin{bytefield}{32}
\bitheader{0-31}\\
\bitbox{16}{Source UDP Port} & \bitbox{16}{Destination UDP Port} \\
\bitbox{16}{Length (header and data)} & \bitbox{16}{Checksum} \\
\bitbox[tlr]{32}[bitheight=3em]{Payload data\\{\scriptsize(variable length, byte aligned)}} \\
\bitbox[lbr]{24}{\hspace{3cm}\raisebox{0.25cm}{$\vdots$}} & \bitbox[lt]{8}{} \\
\end{bytefield}
\end{center}

\begin{exercise}\ \\[-0.5cm]
\begin{itemize}
\item ¿Cuál es el tamaño máximo de payload que se puede enviar en un único \conceptRef{user datagram}{datagrama de usuario}?
\item ¿Es deseable la fragmentación cuando se usa UDP?
\end{itemize}
\end{exercise}

\subsection{Funcionamiento}\label{sec:layer4:udp:operation}

\begin{center}
\includegraphics[width=0.5\linewidth]{udp_socket.pdf}
\end{center}

Los clientes UDP crean primero un \concept{socket} (punto de conexión) y después envían directamente los datagramas de usuario con \inlineCode{sendto}.
No hace falta establecer ni cerrar ninguna conexión.

Los servidores UDP deben hacer primero \concept{bind} al puerto que se va a exponer a los clientes y después
recibir cada datagrama con una llamada individual a la función de lectura (\inlineCode{recvfrom}).

La función bind es puramente interna al servidor: implica cambios en su configuración
(la tabla que relaciona números de puerto y procesos), pero no genera ningún mensaje.

\begin{exercise}\ \\[-0.5cm]
\begin{itemize}
\item ¿Puede funcionar un servidor UDP sin enviar ningún dato por la red?
\item ¿Puede un cliente recibir tráfico UDP sin hacer bind a un puerto?
\item ¿Pueden dos hosts distintos hacer bind al mismo puerto y aun así intercambiar \conceptRef{user datagram}{datagramas de usuario}?
\end{itemize}
\end{exercise}


\begin{exercise}
El siguiente fragmento de código muestra el esqueleto básico de un servidor UDP:

\begin{center}
\showCodeNoOutput{snippets/udpserver.py}
\end{center}

\begin{itemize}
\item Ejecuta el servidor y, en otra terminal, ejecuta \inlineCode{ncat -u localhost 4567} y escribe algo de texto para probar
  el servidor.
\item Implementa el extremo cliente de la comunicación y pruébalo contra el servidor.
\item ¿Qué información importante obtienes si cambias \inlineCode{s.recvfrom} por \inlineCode{s.recvmsg}?
\end{itemize}

\begin{remark}
Este código usa \conceptRef{socket}{sockets} UDP normales, así que no necesita permisos de administrador
ni la capacidad \inlineCode{CAP_NET_RAW}, como sí ocurría en el Ejercicio~\ref{ex:layer2:echo},
que en cambio requería acceso a sockets de capa~2 (\eg\ Ethernet).
\end{remark}
\end{exercise}

\section{TCP -- Transport Control Protocol}\label{sec:layer4:tcp}
\subsection{Formato del paquete}\label{sec:layer4:tcp:packet_format}

Las PDU de TCP se llaman segmentos.
Su cabecera tiene al menos $20$ bytes, pero puede crecer cuando se usan ciertas opciones:

\begin{center}
\begin{bytefield}[bitheight=3em,bitwidth=1.1em]{32}
\bitheader{0-31}\\
\begin{rightwordgroup}{Mandatory\\header part\\{\scriptsize($20$ bytes)}}
\bitbox{16}{Source TCP port} & \bitbox{16}{Destination TCP port} \\
\bitbox{32}{Sequence (seq) number} \\
\bitbox{32}{Acknowledgement (ack) number} \\
\bitheader{0-31}\\
\bitbox{4}{Header\\length/$4$} 
& \bitbox{4}{\otherBase{0}} 
& \bitbox{1}{\rotatebox{90}{\small CWR}} 
& \bitbox{1}{\rotatebox{90}{\small ECE}} 
& \bitbox{1}{\rotatebox{90}{\small URG}} 
& \bitbox{1}{\rotatebox{90}{\small ACK}} 
& \bitbox{1}{\rotatebox{90}{\small PSH}} 
& \bitbox{1}{\rotatebox{90}{\small RST}} 
& \bitbox{1}{\rotatebox{90}{\small SYN}} 
& \bitbox{1}{\rotatebox{90}{\small FIN}} 
& \bitbox{16}{Window size\\{\small (bytes)}}\\
\bitbox{16}{Checksum} & \bitbox{16}{Urgent pointer}
\end{rightwordgroup} \\
\begin{rightwordgroup}[curlyshrinkage=10pt]{Optional\\header part\\{\scriptsize($4$ byte words)}}
\bitbox[tlr]{32}{Options\\{\scriptsize(variable length, optional)}} \\
\bitbox[lbr]{24}{\hspace{3cm}\raisebox{0.25cm}{$\vdots$}} & \bitbox{8}{Padding to\\$32$ bits}
\end{rightwordgroup} \\
\begin{rightwordgroup}[curlyshrinkage=10pt]{Payload}
\bitbox[tlr]{32}[bitheight=3em]{Payload data\\{\scriptsize(variable length, byte aligned)}} \\
\bitbox[lbr]{24}{\hspace{3cm}\raisebox{0.25cm}{$\vdots$}} & \bitbox[lt]{8}{}
\end{rightwordgroup}
\end{bytefield}
\end{center}

Las cabeceras TCP contienen ocho flags (indicadores) de $1$ bit. De ellos, cuatro se usan para funcionalidades centrales de TCP que se tratan en las siguientes
secciones. En el orden en que aparecen:
\begin{itemize}
\item ACK: flag de \conceptRef{ACK}{acuse de recibo} -- ``\textit{mi} número de ack es válido''
\item RST: flag de \conceptRef{RST}{reinicio} -- ``\textit{aborta} la conexión''
\item SYN: flag de \conceptRef{SYN}{sincronización} -- ``estoy abriendo \textit{mi} extremo de la conexión''
\item FIN: flag de \conceptRef{FIN}{finalización} -- ``no enviaré más datos de \textit{payload}''
\end{itemize}

\begin{exercise}
¿Cómo puede saber el extremo receptor el tamaño de las opciones y de los datos de payload?
\end{exercise}

\subsection{Funcionamiento}\label{sec:layer4:tcp:operation}

TCP es un protocolo \concept{stateful} (con estado). Los \conceptRef{socket}{sockets} TCP transitan entre
múltiples estados durante la comunicación. Cada extremo de la comunicación tiene su propio
socket con su propio estado.
% 
A grandes rasgos, las conexiones se pueden clasificar como ``Conectando'',
``Conexión establecida'' o ``Cerrando la conexión''.
% 
Los datos de payload solo están pensados para intercambiarse durante ``Conexión establecida''.

\begin{center}
 \includegraphics[width=0.75\linewidth]{tcp_socket.pdf}
\end{center}

Los sockets TCP pueden estar en uno de $11$ estados. Estos estados, junto con dos contadores de $16$ bits
(\conceptRef{acknowledgement number}{número de confirmación} y \conceptRef{sequence number}{número de secuencia}), hacen posible la \conceptRef{reliability}{fiabilidad} en la comunicación
que se espera de TCP. El significado de estos estados y sus transiciones se explican en los
siguientes capítulos.

\section{TCP: ``Conectando''}\label{sec:layer4:tcp_opening}

\begin{center}
\includegraphics{3_way_handshake.eps}
\end{center}

Cuando se crea un \concept{socket}, empieza como \texttt{CLOSED}.
El extremo que actúa como servidor empieza a escuchar en el puerto configurado para el servicio
y pasa al estado \texttt{LISTEN}. En este estado no se intercambia payload.

Una conexión se establece mediante un \concept{3-way handshake} (negociación inicial) usando los
flags \texttt{SYN} y \texttt{ACK} de la cabecera.
% 
Cada extremo de la comunicación alcanza el estado \texttt{ESTABLISHED} en un momento distinto gracias a
este handshake.
Cuando lo hacen, están listos para intercambiar datos de payload.

\begin{remark}
En TCP, ambos extremos se comunican como peers (pares) iguales; no importa quién inició la comunicación.
Dicho esto, normalmente un extremo realiza una ``apertura activa'' (el cliente)
y el otro una ``apertura pasiva'' (el servidor).
\end{remark}

Cada extremo de la comunicación lleva la cuenta de sus contadores de
\conceptRef{acknowledgement number}{número de confirmación} (\nack) y \conceptRef{sequence number}{número de secuencia} (\nseq).
% 
Cada extremo elige el valor de \nseq\ por separado (y normalmente al azar):
$A,\,B \in [0, 2^{32} - 1]$. 

Una sola vez durante la apertura de la conexión, cada extremo envía exactamente un segmento con el flag \concept{SYN} activado.
Estos segmentos intercambian (sincronizan) sus números de secuencia iniciales $A$ y $B$.

Los segmentos con el flag \concept{ACK} activado (todos salvo el primer segmento del cliente)
tienen su \nack\ fijado a ``el siguiente \nseq\ que se espera del otro extremo''.
Este \nack\ debe leerse como ``he recibido correctamente todos los segmentos correspondientes a
\nack, sin incluirlo''. 


Los números \nack\ y \nseq\ esperados durante el \concept{3-way handshake} son:
\begin{center}
\includegraphics[width=\linewidth]{3_way_handshake_counters.eps}
\end{center}

\begin{exercise}\ \\[-0.5cm]
\begin{itemize}
\item ¿Cómo puede saber un extremo que el otro ha recibido el \nseq\ que ha elegido?
\item ¿Por qué el tercer segmento tiene su \nseq\ fijado a $A+1$ y no a $A$?
\item ¿Qué \nseq\ tendrá el siguiente segmento que envíe el servidor?
\end{itemize}
\end{exercise}


\section{TCP: ``Conexión establecida''}\label{sec:layer4:established}

\subsection{Flujos y números de secuencia}\label{sec:layer4:established:streams_and_sequence_numb}

TCP ofrece la ilusión de un flujo de datos: dos tuberías de datos en los sentidos de entrada y salida
(una conexión \conceptRef{full duplex}{full-duplex}). Estas tuberías son \conceptRef{reliability}{fiables},
\ie, se establecen mecanismos para evitar la pérdida, la duplicación y el reordenamiento de los datos.

\begin{center}
\includegraphics[width=0.75\linewidth]{tcp_two_pipes.pdf}
\end{center}

A cada byte de datos de payload se le asigna un número de secuencia (\nseq) en la cabecera TCP.
El primer byte de datos
de payload tiene \nseq\ $=A+1$, donde $A$ es el número de secuencia inicial establecido durante la apertura
de la conexión (ver Sección~\ref{sec:layer4:tcp_opening}). A cada nuevo byte de datos se le asigna el siguiente
\nseq\ (sin huecos).

\begin{center}
\begin{tabular}{r |c|c|c|c|c|c|c|c|c|c| l}
\textbf{\nseq} & $A$ & $A+1$ & $A+2$ & $A+3$ & $A+4$ & $A+5$ & $A+6$ & $A+7$ & $A+8$ & $A+9$ & $\cdots$ \\[0.2cm]
\textbf{contenido}  & \raisebox{-0.5em}{\rotatebox{90}{SYN}} & 
\begin{minipage}{2.5em}
\centering{\otherBase{0x48} \\ \inlineCode{'H'}}
\end{minipage} &  
\begin{minipage}{2.5em}
\centering{\otherBase{0x69} \\ \inlineCode{'i'}}
\end{minipage} &  
\begin{minipage}{2.5em}
\centering{\otherBase{0x21} \\ \inlineCode{'!'}}
\end{minipage} &  
\begin{minipage}{2.5em}
\centering{\otherBase{0x0a} \\ \inlineCode{'\n'}}
\end{minipage} &  
\begin{minipage}{2.5em}
\centering{\otherBase{0x48} \\ \inlineCode{'H'}}
\end{minipage} &  
\begin{minipage}{2.5em}
\centering{\otherBase{0x6f} \\ \inlineCode{'o'}}
\end{minipage} &  
\begin{minipage}{2.5em}
\centering{\otherBase{0x77} \\ \inlineCode{'w'}}
\end{minipage} &  
\begin{minipage}{2.5em}
\centering{\otherBase{0x20} \\ \inlineCode{' '}}
\end{minipage} &  
\begin{minipage}{2.5em}
\centering{\otherBase{0x61} \\ \inlineCode{'a'}}
\end{minipage} &  
$\cdots$ \\
\end{tabular}
\end{center}

\begin{exercise}
El \nseq\ que viene después de $2^{32} - 1$ es $0$. ¿Por qué?
\end{exercise}

Los datos de payload en TCP siempre se envían junto con la \textit{primera} posición que ocupan en el
flujo y su longitud.
Por ejemplo, a partir de los datos mostrados arriba, se podrían producir dos segmentos:

\begin{enumerate}
\item \nseq\ $= A+1$, len $= 4$, datos $= $ \otherBase{0x48 69 21 0a} (\inlineCode{'Hi!\n'})
\item \nseq\ $= A+5$, len $= 13$, datos $=$ \otherBase{0x48 6f 77 20 61 72 65 20 79 6f 75 3f 0a}
\end{enumerate}

\begin{exercise}
¿Cuál sería el \nseq\ del siguiente segmento con nuevos datos de payload?
\end{exercise}

El extremo receptor de un segmento TCP usa esta información para recomponer
una copia de los datos en el orden correcto, aunque los propios segmentos lleguen desordenados.
%
Estos datos se pasan a la capa superior a través de un
\conceptRef{buffer}{buffer} (memoria intermedia), o directamente si el flag TCP \concept{PSH} (push) está activado.

\begin{center}
\includegraphics[width=\linewidth]{tcp_reassembly.pdf}
\end{center}


\subsection{Fiabilidad y retransmisión}\label{sec:layer4:established:reliability_and_retransmi}
TCP necesita que todos los datos transmitidos sean \conceptRef{ACK}{confirmados} por el otro extremo.
Siempre que se reciben datos en un segmento TCP, el extremo receptor devuelve una confirmación
aunque en ese momento no necesite enviar ningún dato.
El número de confirmación (\nack) de la cabecera TCP se usa para decirle al otro extremo que todos los datos hasta
(sin incluirlo) \nack.

\begin{center}
\includegraphics[width=\linewidth]{tcp_ack.eps}
\end{center}

\begin{remark}
El campo \nseq\ de una cabecera TCP se refiere a los números de secuencia del emisor.
El campo \nack\ se refiere a los números de secuencia del receptor.
\end{remark}

En las confirmaciones TCP se incluyen datos de payload siempre que haya datos disponibles.
A la inversa, el valor más reciente de \nack\ se incluye siempre que se envían datos,
\ie, las confirmaciones hacen \concept{piggyback} (viajan a cuestas) sobre los datos de payload.

Por ejemplo, si Alice envía una petición de $25$ bytes, Bob puede enviar su respuesta de $10$ bytes
y la confirmación de los datos de Alice en el mismo segmento TCP. Alice necesita confirmar
los datos de Bob y, si hace falta, envía enseguida un segmento sin datos. Alice retomará el envío de
datos más adelante cuando lo necesite.

\begin{center}
\includegraphics[width=\linewidth]{tcp_ack_piggyback.eps}
\end{center}

Un segmento TCP que no se confirma tras un cierto periodo de espera se
\conceptRef{retransmission}{retransmite}. Ese tiempo lo decide el sistema operativo de cada extremo
a partir del tiempo medio que tarda en enviarse un segmento y recibirse una respuesta
(el tiempo de ida y vuelta, \concept{RTT}).

Las confirmaciones se envían aunque esa confirmación sea idéntica a otra anterior. Esto es así porque
las confirmaciones también se pueden perder.

\begin{center}
\includegraphics[width=\linewidth]{tcp_ack_retransmission.eps}
\end{center}


\begin{exercise}
Suponiendo que ambos extremos siguen TCP correctamente:
\begin{itemize}
 \item Proporciona un conjunto de valores válidos para (a) a (i).
 \item ¿Cómo cambia la historia si (f) y (h) son idénticos/distintos?
\end{itemize}

\begin{center}
\includegraphics[width=\linewidth]{tcp_ack_exercise.eps}
\end{center}
\end{exercise}

\subsection{Límites de transmisión}\label{sec:layer4:established:transmission_limits}

Los datos que hay que transmitir se suelen producir a ráfagas. Nos gustaría enviar todo lo
posible lo antes posible, pero TCP impone dos límites.

El primer límite, el \conceptRef{MSS}{tamaño máximo de segmento (MSS)}, es la cantidad máxima de datos de payload
por segmento permitida en una conexión TCP.

\begin{center}
\begin{bytefield}[bitheight=3em]{45}
\bitbox{5}{Ethernet header} & 
\bitbox{5}{IP header} & 
\bitbox{5}{TCP header} & 
\bitbox{20}{TCP Payload data}\\[0.25cm]

\bitbox[]{5 }[bitheight=1.5em]{} & 
\bitbox[]{5 }[bitheight=1.5em]{} & 
\bitbox[]{5 }[bitheight=1.5em]{} & 
\bitbox[t]{20}[bitheight=1.5em]{MSS}\\

\bitbox[]{5 }[bitheight=1.5em]{} & 
\bitbox[t]{30}[bitheight=1.5em]{MTU}\\[-0.75cm]
\end{bytefield}
\end{center}

La razón de este límite es evitar la fragmentación, por motivos de rendimiento,
\eg\ por el elevado coste en tiempo y memoria del reensamblado.

Por defecto, el MSS es de $536$~bytes en IPv4,
pero se puede negociar con las opciones de la cabecera TCP durante la fase de apertura de la conexión.
En IPv6, el MSS por defecto es de $1220$~bytes.

\begin{exercise}
Ethernet tiene una MTU de $1500$~bytes. Suponiendo que los datagramas viajan
solo por 2 redes Ethernet:
% 
\begin{itemize}
\item ¿Cuál es el mayor \concept{MSS} que no puede provocar fragmentación, MSS$_\textrm{max}$?
\item ¿Cuál es el overhead para MSS$_\textrm{max}$ y para MSS$_\textrm{max}$ + 1?
\end{itemize}
\end{exercise}

El segundo límite afecta a la cantidad máxima de datos de payload que se pueden enviar
antes de que el otro extremo responda con un nuevo \conceptRef{acknowledgement number}{\nack}:
el llamado tamaño de la \conceptRef{transmission window}{ventana de transmisión}.
% 
Un host no enviará datos (por muchos que haya disponibles o por urgentes que sean) hasta que lleguen esos \nack.
% 
Cuando el ordenador más lento responde, lo hace siempre con el \nack\ más reciente.

\begin{exercise}
¿Cuáles son los tamaños de MSS y de ventana de transmisión más probables en este ejemplo?

\begin{center}
\includegraphics[width=\linewidth]{tcp_ack_window.eps}
\end{center}
\end{exercise}

El valor de la ventana de transmisión se anuncia en todos y cada uno de los segmentos TCP:
este valor es un campo obligatorio de las cabeceras TCP.
% 
A medida que avanza la conexión, cada host puede cambiar el valor que anuncia en función de
su propia carga y necesidades.

\begin{remark}
Cada host anuncia su propio tamaño de ventana.
\end{remark}


Este límite sirve para evitar que los ordenadores rápidos desborden a los más lentos (o sobrecargados):
sin el límite, el ordenador más lento podría necesitar un buffer de tamaño infinitamente creciente para almacenar los
segmentos aún no procesados que le envía el ordenador más rápido. Con este límite,
en el punto de saturación ambos extremos envían segmentos al mismo ritmo, aunque tarden tiempos
muy distintos en producirlos y procesarlos.

\section{TCP: ``Cerrando la conexión''}\label{sec:layer4:tcp_closing}

\subsection{FIN, \nseq\ y \nack}\label{sec:layer4:tcp_closing:fin_nseq_and_nack}

Una vez la conexión está \texttt{ESTABLISHED}, sigue así hasta que uno de los dos extremos decide
cerrarla.

Un nodo anuncia que ha terminado de enviar datos y que quiere cerrar la conexión enviando un
segmento con el flag \concept{FIN} activo en la cabecera. El otro extremo puede seguir
enviando datos, en cuyo caso espera respuestas con valores de \nack\ actualizados. El otro extremo
también puede optar por cerrar la conexión en cuanto ve el primer FIN.
% 
En cualquier caso,
si se envían $k$~bytes de datos y $A$
es el \nseq\ inicial, entonces se espera un segmento con \nack\ $= A+k+2$.

\begin{center}
\begin{tabular}{r |c|c|c|c|c|}
\textbf{\nseq}    & $A$ & $A+1$ & $\cdots$ & $A+k$ & $A+k+1$ \\[0.2cm]
\textbf{contenido}  & \raisebox{-0.5em}{\rotatebox{90}{SYN}} & Datos $\cdots$ & $\cdots$ & $\cdots$ datos & \raisebox{-0.5em}{\rotatebox{90}{FIN}} \\[0.2cm]
\end{tabular}
\end{center}

Fíjate en que el flag FIN ocupa por sí mismo un número de secuencia y debe confirmarse igual que el SYN
(ver Sección~\ref{sec:layer4:tcp_opening}), \eg:

\subsection{Cierre bilateral}\label{sec:layer4:tcp_closing:two_sided_close}

Después de que un extremo cierre su lado de la conexión, el otro puede seguir enviando datos indefinidamente.
Cuando el segundo lado termina de enviar datos, también envía un segmento con el flag \concept{FIN} activado.


\begin{center}
\includegraphics[width=0.9\linewidth]{tcp_close_more_data.eps}
\end{center}

Fíjate en que, si el otro extremo quiere terminar la conexión en cuanto recibe un segmento FIN,
el proceso de cierre se condensa en $3$ segmentos en lugar de $4$, \eg\ 


\begin{center}
\includegraphics[width=\linewidth]{tcp_close_no_more_data.eps}
\end{center}



\begin{remark}
Fíjate en que:
\begin{itemize}
\item Cada extremo envía solo un segmento con el flag FIN, salvo en caso de retransmisión.
\item El segmento con el flag FIN puede contener datos de payload. En ese caso:
  \begin{itemize}
  \item Los datos de payload ocupan números de secuencia \textit{anteriores} al del FIN.
  \item El otro extremo debe confirmar (ACK) tanto los datos de payload como el FIN.
  \end{itemize}
\end{itemize}
\end{remark}


\subsection{Cierre unilateral}\label{sec:layer4:tcp_closing:one_sided_close}

A veces hay que terminar una conexión y no hay tiempo ni interés suficientes
para pasar por el proceso de cierre bilateral de varios pasos.
% 
Puede ser el caso, \eg\ cuando un proceso se cuelga, un puerto TCP o UDP dado está
bloqueado y en otras situaciones de error.

En este escenario, el host que ha detectado la situación anómala simplemente envía un segmento
con su flag \concept{RST} (reset) activado. Los extremos receptores no deben confirmar
el RST; en su lugar, se liberan directamente los recursos de la conexión, \eg\ 

\begin{center}
\includegraphics[width=\linewidth]{tcp_close_reset.eps} 
\end{center}

Seguir el procedimiento de cierre de conexión bilateral o unilateral ayuda a reducir los
recursos de cómputo necesarios para mantener las pilas de red de miles de millones de dispositivos de todo el mundo.

\begin{remark}
Una conexión \textit{no} se puede volver a abrir después de cerrarse. Al flag RST se le podría haber
llamado, con más precisión, ``terminación forzosa''.
\end{remark}

\begin{exercise}
El siguiente diagrama muestra los $11$ estados de TCP y (casi)
todas las transiciones de estado posibles.

\begin{center}
 \includegraphics[width=\linewidth]{tcp_state_transition.pdf}
\end{center}

\begin{itemize}
\item Indica el ciclo de vida del socket de un servidor HTTP como un camino en el grafo.
\item Indica el ciclo de vida del socket del cliente HTTP correspondiente como un camino en el grafo.
\item El estado \inlineCode{CLOSING} no se describe en las secciones anteriores.
  ¿Cuál crees que es su propósito?
\end{itemize}

\begin{remark}
 Este diagrama está adaptado de ``TCP/IP Illustrated, Volume 2: The Implementation,''
de Gary R. Wright y W. Richard Stevens.
\end{remark}


\end{exercise}
